\clearpage
\fancyhead[L]{\small BODE. ASTRONOMICAL YEARBOOK.}
\hypertarget{native-581-lower}{}
\begin{center}\small Göttingische gelehrte Anzeigen. No. 28. Pages 265..277. 18 February 1815.\\[3pt]\rule{.72\linewidth}{.3pt}\end{center}

\emph{Astronomical Yearbook for the Year} 1817 \emph{together with a collection of the latest papers, observations and notices relating to the astronomical sciences. Calculated and edited by} J. E. \textsc{Bode}, \emph{Royal Astronomer and member of the Academy.} Available from the author and distributed on commission by J. E. \textsc{Hitzig}. Berlin 1814. 260 pages in octavo, with a copperplate.

The year 1817 has only two solar eclipses, neither visible in Europe. Among the phenomena that deserve astronomers’ particular attention, we single out a very close conjunction of Venus with Regulus. According to Mr.\ \textsc{Bode}’s calculation, it will take place on 28 September at $14^{\mathrm{u}}$ Berlin time and will therefore be very well placed for observation in eastern Europe. By Mr.\ \textsc{Bode}’s calculation, Venus passes only one minute of arc south of the star and will thus come very close to an actual occultation. Such phenomena deserve to be calculated in advance with the greatest precision. In our view, \emph{actual} occultations of the brighter stars by Venus, particularly when it is nearer inferior conjunction with the Sun, would justify long journeys just as much as transits of Venus across the Sun do. There is also a noteworthy conjunction of Mars with Saturn on 18 April at $7^{\mathrm{u}}$, equally close but invisible in Europe, which would be worth recommending to the observatory at Botanybay for observation.

The \emph{supplements} to this volume of the Yearbook begin with ephemerides for Pallas by Mr.\ \textsc{Nicolai}, for Juno by Mr.\ \textsc{Möbius} and for Vesta by Mr.\ \textsc{Gerling}. These are followed by observations and parabolic elements of the second comet of 1813 by Dr.\ \textsc{Olbers}. --- Astronomical observations made at the Prague Observatory in 1813 by Messrs.\ \textsc{David} and \textsc{Bittner}. The method recommended here for determining azimuths from corresponding horizontal distances from a star at equal altitudes, without using the time, is very good in itself. With an instrument whose horizontal and vertical circles are both sufficiently large and finely graduated, it would also be suitable for very precise determinations. With a theodolite whose altitude circle is graduated only in minutes of arc, or with a circle whose azimuth circle is so graduated, the attainable precision is naturally more limited. With the former instrument one could avoid that drawback by leaving the telescope fixed at the observed altitude, but one would then lose the advantage of repetition. With the circle, small errors in adjusting the transverse level and making the line of sight parallel to the plane of the instrument would also have their full effect on the result. These errors can hardly be avoided altogether, although they do no harm when the instrument is used for its intended purpose. They can be cancelled by repeating the same observations on the following day with the circle in the opposite position, since they must then \SourceLeaf{582}act in the opposite direction as well. --- Some further remarks and additions to the observations on $\eta$ Antinous in the Astronomical Yearbook for 1816, by Prof.\ \textsc{Wurm}. By including ten earlier observations by \textsc{Pigott} that had previously escaped Prof.\ \textsc{Wurm}’s notice, along with a few more of his own, and by applying the method of least squares, this skilled astronomer has now determined the period as 7,17604 days and the epoch as 1800 January 4,504. New tables based on this determination are also given. --- Several valuable astronomical contributions by Prof.\ \textsc{Littrow} in Kazan concern the expansion of the equation of the centre into a series and the theory of epicycles, by which, as the author shows very well, elliptical motions cannot be represented simultaneously in longitude and in radius vector. They also concern a method of determining the circular orbit of a planet from two observations. It is more general and simpler in form than the method developed by the same author in the Yearbook for 1816, but for moderate motions it would give less accuracy in practice because it introduces their cosine, though admittedly the whole problem is at present of theoretical interest only. The same remark applies to determining a straight-line orbit from four observed longitudes (this paper is utterly disfigured by printing errors alone). For the same problem with three complete observations, the author applies the solution borrowed from the \emph{Theoria motus corporum coelestium} to three observations of Vesta (the assumption itself is inadmissible even as an approximation, as shown in the cited work, Art. 131. Mr.\ \textsc{Littrow} proposes it only to make it easier to find a newly discovered celestial body again after a brief interruption, but the usual interpolation methods will already suffice for that purpose). Mr.\ \textsc{Littrow} also gives an indirect solution to the problem of finding the time and latitude from the altitudes of two stars. It is suitable in itself and becomes more convenient when logarithmic differences are used. We must also point out that whenever the hour angles $\alpha$, $\alpha'$ are smaller than $90^{\circ}$, it is better to determine $\frac{1}{2}\alpha$ and $\frac{1}{2}\alpha'$ from their sines by the familiar formulas. We would rank the new method of calculating circum-meridian altitudes below the usual one. The gain in accuracy is too slight, since the error remains of the same order as in the usual method if one keeps only the first term. Likewise, for the last problem, concerning the Earth's spheroid, we prefer the equivalent finite expressions to the series given here. --- Next come astronomical observations at the Vienna Observatory by Messrs.\ \textsc{Triesnecker} and \textsc{Bürg} and at the Cremsmünster Observatory by Mr.\ \textsc{Derfflinger}. --- Observations to determine the geographical position of Port Jackson, calculated from \textsc{Malaspina}’s papers by Mr.\ J. \textsc{Oltmanns}. We are given reason to hope for a complete treatment of all the observations. --- Prof.\ J. W. \textsc{Pfaff} of Nuremberg offers ideas on calculating perturbations after \textsc{Keppler}. They seem to us too superficial to give any result. The form given here for the perturbation equations of the node and inclination (or more precisely for their differential changes) seems to us entirely unproductive. In particular it does not yield the relationships that Mr.\ \textsc{Pfaff} supposes, since, for instance, equations of the same period arise from the terms $Fxx'$ and $Gyy'$. --- Observations at the Wilna Observatory by Prof.\ \textsc{Sniadecky}. --- Prof.\ \textsc{Bürg}’s remarks on determinations of geographical positions in Hungary, Austria and Bavaria contain a defence against a paper in the monatliche Correspondenz, which sought in particular to cast doubt on the azimuth determination made in Vienna. --- On the comet of 1558, by Dr.\ \textsc{Olbers}. Previously only three observations of this comet by Landgrave \textsc{Wilhelm}, on 20, 21 and 23 September, and one by \textsc{Cornelius Gemma} on 20 September were known. Gemma's observation differed greatly from \SourceLeaf{583}the Landgrave’s observation of the same date. By a very fortunate correction of \textsc{Gemma}’s text, reading distabat eo die in place of distabat 20 die, Dr.\ \textsc{Olbers} resolves this contradiction most satisfactorily. \textsc{Gemma}’s observation now applies to 17 August. It makes possible the determination of the comet’s orbit given here, at least to some approximation. --- Some observations in physical astronomy of Saturn, Mars, the Moon, Venus and the Sun, by Dr.\ \textsc{Gruithusen} of Munich, give striking evidence of the excellence of \textsc{Fraunhofer}’s telescopes. --- Also from Prof.\ \textsc{Oltmanns}, a contribution to the determination of the longitude of \textsc{Quito} from several observations by the French academicians, and the calculation of the recently discussed total solar eclipse at Mirabeau in Provence on 3 June 1239. Mr.\ \textsc{Oltmanns} makes a successful application of it to the motion of the Moon’s node. On this evidence the motion used in Mr.\ \textsc{Oltmanns}’s tables needs to be reduced by $1'40''$ over 100 years. By taking two other similar observations into account Mr.\ \textsc{Oltmanns} finds a mean reduction of $1'14''$ in the motion over a century. --- Prof.\ \textsc{Bessel}’s tables for the apparent positions of Polaris are an extraordinarily valuable gift to observing astronomers. It is a pity that the author could not yet use the more recent observations of this star’s right ascension by various astronomers. They agree in showing the need for an increase (at present about 4 seconds of time). We would have liked the pure mean positions for each year to be included as well, to make it easier to compare the determinations of other astronomers. Remarks by Prof.\ \textsc{Bürg} on the revision of his earlier lunar calculations. To improve his lunar tables again Mr.\ \textsc{Bürg} has begun comparing them with the Greenwich observations since 1765. Whatever the results of this laborious work may be, they will certainly be interesting and instructive. --- Observations of Juno and new elements of its orbit, calculation of the next opposition of Pallas and other astronomical notices, by Prof.\ \textsc{Gauss}. --- Some remarks on the prospect of new discoveries in the heavens by means of telescopes. The paper sets out the difficulties that limit further improvements in optical instruments. We must admit, however, that it does not convince us that these obstacles are entirely insurmountable. We must leave open how far the defects highlighted here diminish or even cancel the usefulness of the proposals to make a telescope follow the motion of the stars by clockwork. Nor do we quite understand how it can be claimed on p.\,218 that the continual darting back and forth of the image makes steady observation impossible \emph{during the few seconds of its passage}, when the very effect of the device is to keep the image in the field for a long time. If experience should confirm the force of this objection, it seems to us that mechanics could surely find means, even without teeth and gearing. to give the telescope an extremely gentle and uniform motion. The chief obstacle to extending the power of optical instruments still further probably lies in atmospheric oscillations, which become more noticeable the more the instruments magnify. --- On two variable stars in Hercules, by Dr.\ \textsc{Koch} of Danzig. Unfortunately the Danzig Observatory, which yielded so many useful observations, has also fallen victim to the last siege. The star designated $m$ in Mr.\ \textsc{Bode}’s note, which is still in no catalogue, has already been determined by our Prof.\ \textsc{Harding} and is on sheet XVI of his celestial charts. --- The mean positions of 28 of the most prominent stars in the Pleiades from \textsc{Piazzi}’s new star catalogue, the mean positions of the 36 \textsc{Maskelyne} fundamental stars from the same catalogue, and finally those of Polaris. What is particularly remark\SourceLeaf{584}able about the latter is the parallax in right ascension of $2^{\mathrm{s}}885$ in time that follows from Mr.\ \textsc{Cacciatore}’s observations. It seems too large to us and cannot be reconciled with the numerous observations made by Mr.\ \textsc{von Lindenau} at the Seeberg Observatory. From \textsc{Oriani}’s observed zenith distances we had derived a parallax only half as large (see no. 10 of this journal [p. 574 of this volume]), which admittedly cannot be relied on absolutely either. \textsc{Piazzi} finds the annual proper motion of Polaris in right ascension by comparison

\begin{center}
\small\begin{tabular}{@{}l@{\hspace{8pt}}c@{\hspace{8pt}}r@{}}
with Hevel & $.\;.\;.\;.\;.\;.\;.\;.\;.\;.\;.\;.\;.$ & $+6^{\mathrm{s}}82$\\
Flamstead & $.\;.\;.\;.\;.\;.\;.\;.\;.\;.\;.\;.$ & $+9.03$\\
La Caille & $.\;.\;.\;.\;.\;.\;.\;.\;.\;.\;.\;.$ & $+3.96$\\
Bradley & $.\;.\;.\;.\;.\;.\;.\;.\;.\;.\;.\;.$ & $+1.62$\\
\end{tabular}
\end{center}
\noindent in time. Mr.\ \textsc{von Zach} had found $+3^{\mathrm{s}}177$ by comparing the determinations of \textsc{La Caille} and \textsc{Mathieu}. It would be highly desirable to have this important subject carefully investigated, though all the original observations would first have to be reduced in a uniform way. --- Astronomical observations at the Royal Observatory in Berlin in 1813. --- Observations at the Dorpat Observatory by Prof.\ \textsc{Struve} in 1812, 1813 and 1814. --- Several further observations at Cremsmünster. --- Further notices on the double star 61 Cygni from \textsc{Piazzi}’s Libro sesto del reale osservatorio (where the remarkable proper motion of this double star was first reported) and from the same astronomer’s new large star catalogue. --- A demonstration that none of eight fixed stars missing from the sky was Ceres, Pallas, Juno or Vesta. We note here that all these stars except the one in the Northern Crown had already been discussed in volume nine of the monatliche Correspondenz (pp. 153. 151. 241. 155. 243. 241. 246), and that in particular all the stars of \textsc{Mayer} mentioned here had already been satisfactorily accounted for from the original papers. ---

Among the miscellaneous remarks and notices appended as usual at the end of the Yearbook, we single out just one. It is a passage from a letter of 6 May 1814 from Dr.\ \textsc{Reche} in Mühlheim on the Rhine. According to it, during the blockade of Strasbourg Prof.\ \textsc{Kramp} found a new solution to the problem of determining a comet’s orbit from three observations. He found that all the cometary orbits he calculated that had previously been regarded as parabolas are \emph{quite certainly and definitely} hyperbolas. Such an announcement must be very striking, and Mr.\ \textsc{Kramp}’s work has since appeared in the July issue of the mathematical journal that has been published for some years in Nimes under the title Annales de mathématiques pures et appliquées. This is therefore the place to add an assessment of that paper to the foregoing review.

Two points matter here. One is the value of the method itself. The other is how it is applied to comets actually observed and how reliable the resulting findings are. We will first make a few remarks on the latter. This paper discusses the application to only one comet, namely that of 1781. From the observations of 14, 19. and 25. November, Mr.\ \textsc{Kramp} derives a hyperbola with eccentricity $= 4,580652$ and perihelion distance $1,048364$. Anyone with even a little experience in calculating cometary orbits would find the short interval enough to judge the reliability of this result. We must add, however, that the resulting elements are far from being tested against other observations. They are not even compared with the three observations on which they are based. Such a check is at least made very difficult for anyone who might wish to carry it out independently, since Mr.\ \textsc{Kramp} has not even carried the calculation of the elements through to the \SourceLeaf{585}end, so that the longitude of perihelion and time of passage are missing. If one finally notes that the Earth’s orbit is treated as a circle throughout the calculation, one will probably think it needless to investigate separately whether there might also be errors in the numerical calculation itself. The \emph{considerable} difference between the value of the quantity
\[
P\!\cdot\!Q'' - P''Q = 0.4614411
\]
and the sum of the two quantities
\[
\begin{aligned}
PQ' - P'Q &= 0.2133547\\
P'Q'' - P''Q' &= 0.2495918
\end{aligned}
\]
which, as a consequence of the method, should actually be \emph{exactly} equal to the first quantity, seems to point to such errors. (Mr.\ \textsc{Kramp} takes a different view and calls this considerable discrepancy an egalité presque rigoureuse, but what is the point of calculating to seven decimal places if one cannot even vouch for the third?) We say \emph{seems}, because this difference might possibly be due to the omissions on p. 13.

We now turn to the method itself. It has already been pointed out that Mr.\ \textsc{Kramp} (quite needlessly, in fact) entirely neglects the eccentricity of the Earth’s orbit. The essential point of the method can be stated in a few words. Mr.\ \textsc{Kramp} denotes the semi-major axis of the celestial body’s orbit by $b$. He denotes the coordinates of its three positions in the orbit, referred to the line of nodes as the axis of abscissas, by $bP$, $bP'$, $bP''$, and $bQ$, $bQ'$, $bQ''$. As is well known, these six quantities are easily determined from the observational data, the inclination of the orbit $\beta$ and the longitude of the ascending node $\delta$. On the other hand, by combining the formulas relating to \textsc{Kepler}’s laws and setting the differences of the eccentric anomalies equal to their sines, Mr.\ \textsc{Kramp} finds the three quantities
\[
\begin{gathered}
PQ' - P'Q\\
P'Q'' - P''Q'\\
PQ'' - P''Q
\end{gathered}
\]
to be proportional respectively to the intervals $t$, $t'$, $t+t'$, giving the equation
\[
0 = PQ' - P'Q + P'Q'' - P''Q' + P''Q - PQ''
\]
which Mr.\ \textsc{Kramp} calls an équation essentielle et très remarquable par sa simplicité. --- In the two equations
\[
\begin{aligned}
(t+t')(PQ'-P'Q) &= t(PQ''-P''Q)\\
(t+t')(P'Q''-P''Q') &= t'(PQ''-P''Q)
\end{aligned}
\]
Mr.\ \textsc{Kramp} substitutes for $P$, $P'$, $P''$, $Q$, $Q'$, $Q''$ their values in terms of $\beta$ and $\delta$ (where $b$ evidently drops out). He obtains two equations in $\beta$ and $\delta$, from which, by omitting a couple of terms and eliminating $\beta$, he finally derives a cubic equation for $\operatorname{tang}\delta$. This determines $\delta$. Then one of the two equations just mentioned (or both) determines $\beta$. The next quantities to be determined are $bP$, $bP'$, $bP''$, $bQ$, $bQ'$, $bQ''$. From four of these quantities one obtains the remaining elements of the conic section by methods we omit here for brevity.

It will be very easy to prove that this method is wholly inadmissible. Here we see the laws of motion replaced by the proportionality to the time intervals of the three quantities mentioned several times. The whole solution rests solely on that proportionality, or the two equations \SourceLeaf{586}equivalent to it, and on the proposition that the three positions lie in one plane. The three quantities
\[
bb(PQ'-P'Q), \quad bb(P'Q''-P''Q'), \quad bb(PQ''-P''Q)
\]
are, as one will immediately notice, precisely the three quantities denoted by $n''$, $n$, $n'$ in the Theoria motus Corporum Coelestium, Art. 112. But it is proved there in Art. 114 that the distances from the Earth are determined \emph{linearly} as soon as the ratio of these three quantities is taken as known. It is therefore an error for Mr.\ \textsc{Kramp} to reduce to a cubic equation a problem whose \emph{rigorous} solution requires only an equation of the first degree, even after neglecting some terms that, if included in Mr.\ \textsc{Kramp}’s treatment, would have led to a biquadratic equation. This, however, is even less important than the second point. Art. 131 of the cited work already proves in detail that it is inadmissible to set those three quantities $n''$, $n$, $n'$ proportional to the intervals $t$, $t'$, $t+t'$ \emph{simultaneously}, since the result cannot be accepted even as an approximation.

If, as it almost seems, Mr.\ \textsc{Kramp} was unfamiliar with the Theoria motus corporum coelestium, which appeared five years before the blockade of Strasbourg, he could still easily have satisfied himself that the true spirit of his method rests on a principle whose unsoundness has long been known. For, as a little attention immediately shows, the three quantities mentioned above are simply twice the areas of the triangles between the first and second, second and third, and first and third radius vectors. To set these three quantities proportional to the three time intervals \emph{simultaneously} is thus simply to assume the three positions of the comet to lie on one straight line, the two segments of which are proportional to the time intervals. And so Mr.\ \textsc{Kramp}’s so-called équation essentielle is simply the condition equation for the three positions to lie on one straight line. It has long been known that in the comet problem this assumption must lead to a solution requiring only equations of the first degree and that the assumption itself is inadmissible even as an approximation. A proof as thorough as it is lucid can be found in particular in \textsc{Olbers}’s excellent work, which Mr.\ \textsc{Kramp} himself cites in the present paper.
